\section{Starting Point: Fundamental Formalism}

We formally define the Hilbert space $\mathcal{H} = L^2(\mathbb{R}^3)$ and the
dynamical structure of the free particle. The following definitions establish
the algebraic and spectral foundations required for the angular-momentum
decomposition.

\subsection{Dynamics and Unitary Evolution}

The dynamics are governed by the self-adjoint Hamiltonian operator $\hat{H} =
\frac{\hat{\mathbf{p}}^2}{2m}$. The state vector $\ket{\psi(t)}$ evolves
deterministically from an initial state $\ket{\psi_0}$ via the strongly
continuous unitary group $\hat{U}(t)$:
\begin{equation}
  \ket{\psi(t)} = \hat{U}(t)\ket{\psi_0} =
  \exp\left(-\frac{i}{\hbar}\hat{H}t\right)\ket{\psi_0}.
\end{equation}
This abstract evolution implies the Schrödinger equation $i\hbar \partial_t
\ket{\psi} = \hat{H}\ket{\psi}$ via Stone's Theorem.

\subsection{Continuous Representations}

We utilize two continuous bases, position $\{\ket{\mathbf{r}}\}$ and linear
momentum $\{\ket{\mathbf{p}}\}$, normalized to the Dirac distribution.

\subsubsection{Basis Definitions and Orthogonality}
The basis vectors satisfy the eigenvalue equations and completeness relations
(Resolution of Identity $\hat{\mathbb{I}}$):
\begin{align}
  \hat{\mathbf{r}}\ket{\mathbf{r}} &= \mathbf{r}\ket{\mathbf{r}}, \quad &
  \int_{\mathbb{R}^3} \mathrm{d}^3r \, \ket{\mathbf{r}}\bra{\mathbf{r}} &=
  \hat{\mathbb{I}}, \\
  \hat{\mathbf{p}}\ket{\mathbf{p}} &= \mathbf{p}\ket{\mathbf{p}}, \quad &
  \int_{\mathbb{R}^3} \mathrm{d}^3p \, \ket{\mathbf{p}}\bra{\mathbf{p}} &=
  \hat{\mathbb{I}}.
\end{align}
The inner product defines the transformation kernel (Fourier kernel) between the
canonically conjugate variables:
\begin{equation}
  \braket{\mathbf{r}}{\mathbf{p}} = \frac{1}{(2\pi\hbar)^{3/2}} \exp\left(
  \frac{i}{\hbar} \mathbf{p}\cdot\mathbf{r} \right).
\end{equation}

\subsubsection{Wavefunction Representations}
The projections of the state vector yield the coordinate and momentum space
wavefunctions:
\begin{itemize}
  \item \textbf{Coordinate Space:} $\psi(\mathbf{r}, t) \equiv
    \braket{\mathbf{r}}{\psi(t)}$. The momentum operator acts as the generator
    of spatial translations:
    \begin{equation}
      \bra{\mathbf{r}}\hat{\mathbf{p}}\ket{\psi} = -i\hbar \nabla
      \psi(\mathbf{r}).
    \end{equation}
  \item \textbf{Momentum Space:} $\tilde{\psi}(\mathbf{p}, t) \equiv
    \braket{\mathbf{p}}{\psi(t)}$. Since $\hat{H}$ is diagonal in this basis,
    the time evolution is strictly a phase modulation:
    \begin{equation}
      \tilde{\psi}(\mathbf{p}, t) = \exp\left( -\frac{i}{\hbar}
      \frac{p^2}{2m} t \right) \tilde{\psi}(\mathbf{p}, 0).
    \end{equation}
\end{itemize}

\subsection{Spectral Decomposition and Density of States}

For the free particle, the energy $E$ is degenerate. The transition from
momentum integration to energy integration requires the spectral density
function. Using the dispersion relation $E_p = p^2/2m$:
\begin{equation}
  \int \mathrm{d}^3p = \int_0^\infty p^2 \mathrm{d}p \int_{S^2}
  \mathrm{d}\Omega.
\end{equation}
Changing variables from $p$ to $E$ introduces the density of states factor via
the identity:
\begin{equation}
  \mathrm{d}p = \frac{\mathrm{d}p}{\mathrm{d}E} \mathrm{d}E = \frac{m}{p}
  \mathrm{d}E = \sqrt{\frac{m}{2E}} \mathrm{d}E.
\end{equation}
This transformation is critical for isolating the radial (energy) dependence
from the angular dependence in the final construction.

\subsection{Symmetry Generators}

The continuous symmetries of the system are generated by Hermitian operators,
implying specific conservation laws via Ehrenfest's Theorem,
$\frac{\mathrm{d}}{\mathrm{d}t}\braket{\hat{A}} = \frac{1}{i\hbar}\braket{[\hat{A},
\hat{H}]}$.

\begin{enumerate}
  \item \textbf{Spatial Translations:} Generated by $\hat{\mathbf{p}}$. Since
    $[\hat{\mathbf{p}}, \hat{H}] = 0$, linear momentum is conserved.
    \begin{equation}
      \hat{T}(\mathbf{a}) = \exp\left( -\frac{i}{\hbar} \mathbf{a} \cdot
      \hat{\mathbf{p}} \right).
    \end{equation}
  \item \textbf{Galilean Boosts:} Generated by $\hat{\mathbf{r}}$ (up to mass
    scaling). This shifts the momentum spectrum:
    \begin{equation}
      \hat{T}_{\mathbf{p}}(\mathbf{q}) = \exp\left( \frac{i}{\hbar}
      \mathbf{q} \cdot \hat{\mathbf{r}} \right).
    \end{equation}
\end{enumerate}
